\scanpage{43}
D'autre part, on a
\[
 R(D/U,k)\hookrightarrow R(D,k)
\]
et l'image de $R(D/U,k)$ dans $R(D,k)$ est précisément \unsure{formé|formée} des classes des $E$ \unsure{satisfaisant|satisfaisants} la condition du début. Pour un tel $E$, on a donc une opération $\mathord{\text{\tikz[x=.66em,y=.68em,baseline=0pt]{\draw[line join=round] (.80,.95)--(.12,.24)--(.83,.34)--(.14,-.37);}}}$ sur la \uline{classe} de $E$. Donc $\operatorname{cl}(E)^{\mathord{\text{\tikz[x=.66em,y=.68em,baseline=0pt]{\draw[line join=round] (.80,.95)--(.12,.24)--(.83,.34)--(.14,-.37);}}}}$ est défini dès qu'il existe un ss-gpe \deleted{\unsure{d'ind|d'indice}} $U$ d'indice fini dans $I$ tel que $\deleted{D/U}$ $U$ opère sur $E$ de façon unipotente. / Si on prend des classes de repr. \uline{continues} ($k$ étant un corps top.) on \deleted{\unsure{parler|parlera|prendra}} \inserted{\unsure{obtient|obtiendra}} de $=$ $\operatorname{cl}(E)^{\mathord{\text{\tikz[x=.66em,y=.68em,baseline=0pt]{\draw[line join=round] (.80,.95)--(.12,.24)--(.83,.34)--(.14,-.37);}}}}$ comme classe de repr. \uline{continues} de $D/I$. Compatibilité avec ext du corps de base.

\deleted{Supposons que $D/I\simeq\mathbb{Z}$, \unsure{et notons|et disons|et soit|et donc toute} qu'une classe de repr. sur $D/I$ est \unsure{donné|donnée|déterminé|déterminée} par \unsure{sa} série car. $-$ $L_{E^{\mathord{\text{\tikz[x=.66em,y=.68em,baseline=0pt]{\draw[line join=round] (.80,.95)--(.12,.24)--(.83,.34)--(.14,-.37);}}}}}(t)$. \unsure{On|Ceci|Comme} \unsure{aura|arrive} alors}
% archive sheet number: 42
